\begin{satz}
{Obere Schranke des Erwartungswerts von abs.
\textsl{N}(0,1)--Ord.st.}
{ObereAbschaetzungDesErwartungswertsDerKtenAbsolutenN01Ordnungsstatistik}
Seien $k, n \in \N$ mit $k \le \frac n 2$, $\gamma:\Omega\rightarrow \R^n$
unabhängig und $N(0,1)$--verteilt und $u: [0, \infty[ \rightarrow \R,
t \mapsto \Prim(\abs{\gamma_1}> t)$.\\
Dann gilt:
\begin{enumerate}[(1)]
	\item\label
	{ObereAbschaetzungDesErwartungswertsDerKtenAbsolutenN01Ordnungsstatistik1}
	$\forall\, x \in \R_{>0}: u(x) \le \frac k n
	\folgt \E\left(\left(\gamma_k^\ast\right)^2\right) \le x^2 + 4$.
	\item\label
	{ObereAbschaetzungDesErwartungswertsDerKtenAbsolutenN01Ordnungsstatistik2}
	$\E\left(\left(\gamma_k^\ast\right)^2\right) 
	\le 2 \log\left(\frac n k \right) + 4 
	\le \left(2+\frac 4{\log(2)} \right)\log\left(\frac n k \right)$.
	\item\label
	{ObereAbschaetzungDesErwartungswertsDerKtenAbsolutenN01Ordnungsstatistik3}
	$\E\left(\left(\gamma_k^\ast\right)^2 \right)
	\le 2 \log\left(\frac {2n} k \right) + 4 
	\le \left(2+\frac 4{\log(4)} \right)\log\left(\frac {2n} k \right)$.
	\item\label
	{ObereAbschaetzungDesErwartungswertsDerKtenAbsolutenN01Ordnungsstatistik4}
	$\E\left(\gamma_k^\ast\right) \le \sqrt{2+\frac 4{\log(2)}}
	\sqrt{\log\left(\frac n k\right)}$.
	\item\label
	{ObereAbschaetzungDesErwartungswertsDerKtenAbsolutenN01Ordnungsstatistik5}
	$\E\left(\gamma_k^\ast\right) \le \sqrt{2+\frac 4{\log(4)}}
	\sqrt{\log\left(\frac {2n} k\right)}$.
	\item\label
	{ObereAbschaetzungDesErwartungswertsDerKtenAbsolutenN01Ordnungsstatistik6}
	Ist $k \le \frac n {\exp(1)}$, so ist $\E\left(\norm\cdot_{2,k} \circ 
	\left(\gamma_j^\ast\right)_{j\in \haken k}   \right)
	\le \sqrt 2 \sqrt{2+\frac 4{\log(2)}}\sqrt{k\log\left(\frac{n}k\right)}$.
	\item\label
	{ObereAbschaetzungDesErwartungswertsDerKtenAbsolutenN01Ordnungsstatistik7}
	$\E\left(\norm\cdot_{2,k} \circ 
	\left(\gamma_j^\ast\right)_{j\in \haken k}   \right)
	\le \sqrt 2\sqrt{2+\frac 4{\log(4)}}\sqrt{k\log\left(\frac{2n}k\right)}$.
\end{enumerate}
\end{satz}
\begin{proof}
(1) Zunächst gilt offenbar:
\begin{align}
\forall\, t \in [0, \infty[: \left(\gamma_k^\ast\right)^2 > t \iff
\gamma_k^\ast > \sqrt t \text.\label{oadedkanoeq:1}
\end{align}
Sei $x \in \R_{>0}$ mit $u(x) \le\frac k n $.
Wegen $\haken{n-k}_0 \subseteq \haken {n-1}_0$ gilt nach
\refclaim{AbschaetzungDerUnvollstaendigenBetafunktion}
{AbschaetzungDerUnvollstaendigenBetafunktion2}:
\begin{align}\forall\, l \in \haken {n-k}_0:
\int_x^\infty (1-u(t))^l(u(t))^{n-l}tdt
\le \int_0^{u(x)} (1-t)^lt^{n-l-1}dt \label{oadedkanoeq:2}
\end{align}
und wegen $u(x) \le \frac k n \le \frac k {n-1}$ nach
\refclaim{AbschaetzungDerUnvollstaendigenBetafunktion}
{AbschaetzungDerUnvollstaendigenBetafunktion1}:
\begin{align}\forall\, l \in \haken{n-k-1}_0:
\int_0^{u(x)} t^{n-l-1}(1-t)^ldt
\le (u(x))^{n-l}(1-u(x))^l \text.\label{oadedkanoeq:3}
\end{align}
Seien $B, \Gamma$ wie in~\ref{ElementareEigenschaftenDerGammafunktion}.
Wir fassen $\gamma_k^\ast$ als $k$--te Ordnungsstatistik von
$\pfeil {\abs\cdot} n \circ \gamma$ auf.\\
Dann gilt mit \refclaim{EigenschaftenVonOrdnungsstatistiken}
{EigenschaftenVonOrdnungsstatistiken2}:
\begin{align*}
& \int_{x^2}^\infty \Prim\left(\gamma_k^\ast > \sqrt t\right)dt\\
\overset{\refclaim{EigenschaftenVonOrdnungsstatistiken}
{EigenschaftenVonOrdnungsstatistiken3}}= &\,
\int_{x^2}^\infty \sum_{l=0}^{n-k}\binom n l 
\Prim\left(\abs{\gamma_1} \le\sqrt t\right)^l
\left(1-\Prim\left(\abs{\gamma_1}\le \sqrt t\right)\right)^{n-l}dt\\
=&\, \sum_{l=0}^{n-k}\binom n l \int_{x^2}^\infty 
\left(u\left(\sqrt t\right)\right)^{n-l}\left(1-u\left(\sqrt t\right)\right)^l
dt\\
\overset{x>0}=&\, \sum_{l=0}^{n-k}\binom n l \int_{x^2}^\infty 
\left(u\left(\sqrt t\right)\right)^{n-l}\left(1-u\left(\sqrt t\right)\right)^l
2\sqrt t \frac 1 {2\sqrt t}dt\\
\overset{\substack{\text{Subst.}\\t \mapsto \sqrt t}}= &\,
2\sum_{l=0}^{n-k}\binom n l \int_x^\infty (u(t))^{n-l}(1-u(t))^l tdt\\
\overset{\eqref{oadedkanoeq:2}}\le &\,
2\sum_{l=0}^{n-k}\binom n l \int_0^{u(x)} t^{n-l-1}(1-t)^ldt\\
\overset{\eqref{oadedkanoeq:3}}\le &\,
2\binom n {n-k}\int_0^{u(x)}t^{k-1}(1-t)^{n-k}dt +
2\underbrace{
\sum_{l=0}^{n-k-1}\binom n l (u(x))^{n-l}(1-u(x))^l}_
{\overset{\text{Bin. Form.}}\le 1}\\
\le &\, 2 \binom n{n-k}B(k, n-k+1) + 2 
\overset{\refclaim{ElementareEigenschaftenDerGammafunktion}
{ElementareEigenschaftenDerGammafunktion4}}=
2 \frac{n!}{(n-k)!k!}\frac{\Gamma(k)\Gamma(n-k+1)}{\Gamma(n+1)} + 2\\
\overset{\refclaim{ElementareEigenschaftenDerGammafunktion}
{ElementareEigenschaftenDerGammafunktion3}}=&\,
2 \frac{n!}{(n-k)!k!} \frac{(k-1)!(n-k)!}{n!} + 2
= \frac 2 k + 2 \le 4\text.
\end{align*}
Somit gilt wegen $\gamma_k^\ast \ge 0$:
\begin{align*}
\E\left(\left(\gamma_k^\ast\right)^2 \right) 
&= \int_0^\infty \Prim\left(\left(\gamma_k^\ast\right)^2> t\right)dt
\overset{\eqref{oadedkanoeq:1}}=
\int_0^\infty \Prim\left(\gamma_k^\ast> \sqrt t\right)dt\\
&= \int_0^{x^2}\underbrace{\Prim\left(\gamma_k^\ast> \sqrt t\right)}_{\le 1}dt
+ \int_{x^2}^\infty \Prim\left(\gamma_k^\ast> \sqrt t\right)dt
\le x^2 + 4\text.
\end{align*}
(2) Sei $x \coloneqq  u\inv\left(\frac k n \right)$.\\
Dann ist wegen $3 < \pi < 4$ und der Standardabschätzung für $\exp$
\[\exp\left(\frac 1\pi\right)\le\frac 1{1-\frac1\pi}=\frac\pi{\pi-1}< 2\text,\]
also auch $\exp\left(-\frac1\pi \right) \ge \frac 1 2 \ge \frac k n$ und somit
nach \refclaim{EigenschaftenDerKomplementaerenVerteilungsfunktionVonN01}
{EigenschaftenDerKomplementaerenVerteilungsfunktionVonN017}:
$x \le \sqrt{2 \log\left(\frac n k \right)}$.\\
Somit
\begin{align*}
\E\left(\left(\gamma_k^\ast\right)^2 \right) \overset{(1)} \le
2 \log\left(\nicefrac n k \right) + 4 
= \left(2 + \nicefrac 4 {\log\left(\frac n k \right)} \right)
\log\left(\nicefrac n k \right)
 \overset{\frac n k \ge 2}\le
\left(2 + \nicefrac 4 {\log(2)} \right)
\log\left(\nicefrac n k \right)\text.
\end{align*}
(3) Es gilt:
\begin{align*}
\E\left(\left(\gamma_k^\ast\right)^2 \right) &\overset{(2)} \le
2 \log\left(\frac n k \right) + 4
\le 2 \log\left(\frac {2n} k \right) + 4 = \left(2 + \frac 4 {\log\left(\frac {2n} k \right)} \right)
\log\left(\frac {2n} k \right)\\
&\overset{\frac{2n}k \ge 4}\le
\left(2 + \frac 4 {\log(4)} \right)\log\left(\frac {2n} k \right)\text.
\end{align*}
(4)/(5) Gemäß Jensenungleichung ist
$\E\left(\gamma_k^\ast \right)= 
\E\left(\sqrt\cdot \circ \left(\gamma_k^\ast \right)^2 \right)
\le \sqrt{\E\left(\left(\gamma_k^\ast\right)^2 \right)}$, also
folgt (4) aus (2) und (5) aus (3).\\
(6) Sei $k \le \frac n {\exp(1)}$. Dann gilt:
\begin{align*}
& \E\left(\norm\cdot_{2,k} \circ 
	\left(\gamma_j^\ast\right)_{j\in \haken k}   \right)
= \E\left(\sqrt\cdot \circ \sum_{j=1}^k \left(\gamma_j^\ast\right)^2 \right)
\le \sqrt{\sum_{j=1}^k \E\left(\left(\gamma_j^\ast\right)^2 \right)}\\
\overset{(2)}\le&\, \sqrt{\sum_{j=1}^k\left(2+\frac 4{\log(2)} \right)
\log\left(\frac n j\right)}
=\sqrt{2+\frac 4{\log(2)}}\sqrt{\sum_{j=0}^k\log\left(\frac n j\right)}\\
\overset{\ref{SummenVonLogarithmen}}\le&\,
\sqrt{2+\frac 4{\log(2)}}\sqrt{2 k \log\left(\frac n k \right)}
= \sqrt 2 \sqrt{2+\frac 4{\log(2)}}\sqrt{k \log \left(\frac n k \right)}\text.
\end{align*}
(7) Wegen $k \le \frac n 2 = \frac{2n} 4 \le \frac{2n}{\exp(1)}$ folgt diese
Aussage analog zu (6), nur dass (3) statt (2) angewandt wird.
\end{proof}